\documentclass[12pt, a4paper]{article}

\usepackage[margin=2.5cm]{geometry}
\usepackage{booktabs}
\usepackage{graphicx}
\usepackage{hyperref}
\usepackage{amsmath}
\usepackage{setspace}
\usepackage{caption}
\usepackage{float}
\usepackage{array}
\usepackage{parskip}
\usepackage{fancyhdr}
\usepackage{microtype}

\pagestyle{fancy}
\fancyhf{}
\rhead{Parking Slot Availability Prediction}
\lhead{Abhijith Mohanan}
\rfoot{Page \thepage}

\hypersetup{
    colorlinks=true,
    linkcolor=blue,
    urlcolor=blue,
    citecolor=blue
}

\onehalfspacing

\begin{document}

% ── Title Page ───────────────────────────────────────────────────
\begin{titlepage}
    \centering
    \vspace*{2cm}
    {\LARGE \textbf{Parking Slot Availability Prediction}}\\[0.5cm]
    {\large \textbf{San Diego Parking Meter Data --- Binary Classification}}\\[2cm]
    {\large \textbf{Abhijith Mohanan}}\\[0.3cm]
    {\normalsize MSc Data Science \& Intelligent Analytics}\\[0.1cm]
    {\normalsize FH Kufstein Tirol --- University of Applied Sciences}\\[2cm]
    {\normalsize \today}\\[3cm]
    \vfill
    {\small
    Datasets: City of San Diego Open Data Portal\\
    \url{https://data.sandiego.gov}
    }
\end{titlepage}

% ── Abstract ─────────────────────────────────────────────────────
\begin{abstract}
Urban parking availability is a growing challenge in cities worldwide.
Drivers waste significant time and fuel searching for available parking,
contributing to traffic congestion and emissions. This project applies
machine learning to predict whether a parking meter transaction in San
Diego will result in a Long Stay (more than 60 minutes) or a Short Stay
(60 minutes or less), using real historical data from the City of San
Diego Open Data Portal. Two datasets are combined --- parking meter
transactions from 2020 (1,953,052 records) and active meter location
data (4,060 meters) --- resulting in a feature-engineered dataset of
1,474,238 records with 19 predictive features. Two classification models
are trained and compared: Random Forest and XGBoost. Both models achieve
over 86\% accuracy on the test set, with XGBoost performing marginally
better overall (F1 Score: 0.8596 vs 0.8590). The results show that
transaction amount and payment method are the strongest predictors of
parking duration, with time-of-day and zone features also contributing
significantly.
\end{abstract}

\newpage
\tableofcontents
\newpage

% ── 1. Introduction ───────────────────────────────────────────────
\section{Introduction}

Finding parking in a busy city is a daily frustration for millions of
drivers. Beyond the inconvenience, time spent searching for available
spots contributes measurably to urban traffic congestion and vehicle
emissions. Smart parking systems --- those equipped with sensors, meters,
and data logging capabilities --- offer an opportunity to turn this
challenge into a data problem.

The City of San Diego operates a large network of smart parking meters
across downtown, uptown, and surrounding zones. Each payment transaction
is automatically recorded with details including the start time, expiry
time, meter zone, area, and payment method. This rich transactional
history provides a natural basis for predicting parking behavior.

\subsection*{Objectives}
This project has three main objectives:
\begin{enumerate}
    \item Combine and clean two real-world parking datasets to build
    a reliable, feature-rich dataset for modeling.
    \item Train and compare two classification models --- Random Forest
    and XGBoost --- to predict whether a parking transaction will be a
    Long Stay (over 60 minutes) or Short Stay.
    \item Evaluate model performance using standard classification metrics
    and identify the most predictive features.
\end{enumerate}

The prediction task is framed as binary classification rather than
regression because the practical value lies in giving drivers and
planners a clear, actionable signal --- not an exact duration estimate.

% ── 2. Related Work ───────────────────────────────────────────────
\section{Related Work}

Parking availability prediction has been approached from several angles
in the literature. Sensor-based approaches using occupancy sensors have
shown strong results in controlled environments \cite{zheng2015}.
However, these require expensive hardware infrastructure.
Transaction-based approaches, like the one used in this project, leverage
existing payment data and are more scalable \cite{shoup2006}.

Machine learning methods have been widely applied to parking prediction
tasks. Ji et al. \cite{ji2014} demonstrated that time-of-day and
day-of-week features consistently rank among the most predictive variables
across multiple urban datasets. Ensemble methods such as Random Forest and
Gradient Boosting have shown particularly strong performance on tabular
urban data \cite{chen2016}.

More recently, studies using the UCI Birmingham Parking Dataset
\cite{uci2016} highlighted that occupancy patterns in city-centre zones
differ significantly from suburban zones --- a pattern also observed in
the San Diego data, where the City and Downtown zones consistently show
higher long-stay rates than Uptown or Mid-City.

% ── 3. Dataset Description ────────────────────────────────────────
\section{Dataset Description}

\subsection{Data Sources}
Both datasets are sourced from the City of San Diego Open Data Portal
(\url{https://data.sandiego.gov}), which provides freely accessible
government data under an open license.

\subsection{Parking Meter Transactions (2020)}
This dataset contains every parking payment recorded in San Diego
during 2020, from February to December. It contains 1,953,052 rows
and 6 columns, as described in Table~\ref{tab:transactions}.

\begin{table}[H]
\centering
\caption{Parking Meter Transactions Dataset Columns}
\label{tab:transactions}
\begin{tabular}{@{}lll@{}}
\toprule
\textbf{Column} & \textbf{Type} & \textbf{Description} \\
\midrule
pole\_id            & String   & Unique meter identifier \\
meter\_type         & String   & Type of meter (all SS in 2020) \\
date\_trans\_start  & Datetime & Transaction start timestamp \\
date\_meter\_expire & Datetime & Meter expiry timestamp \\
trans\_amt          & Integer  & Amount paid (in cents) \\
pay\_method         & String   & CASH, CREDIT CARD, or SMART CARD \\
\bottomrule
\end{tabular}
\end{table}

Payment method breakdown: Cash (1,058,818), Credit Card (892,069),
Smart Card (2,165). Date range: 24 February 2020 to 31 December 2020.
Unique meters in transactions: 4,592.

\subsection{Parking Meter Locations (Active)}
This dataset contains static information about all active parking meters
in San Diego. It contains 4,060 rows and 16 columns, with key fields
described in Table~\ref{tab:locations}.

\begin{table}[H]
\centering
\caption{Parking Meter Locations Dataset --- Key Columns}
\label{tab:locations}
\begin{tabular}{@{}lll@{}}
\toprule
\textbf{Column} & \textbf{Type} & \textbf{Description} \\
\midrule
pole                & String  & Meter ID (join key) \\
zone                & String  & Parking zone \\
area                & String  & Sub-area within zone \\
latitude            & Float   & GPS latitude \\
longitude           & Float   & GPS longitude \\
time\_limit         & String  & Maximum parking duration \\
days\_in\_operation & String  & Operating days \\
price               & String  & Hourly rate \\
\bottomrule
\end{tabular}
\end{table}

Zone distribution: Downtown (2,404), Uptown (1,090), Mid-City (348),
City (90), Balboa Park (60), Pacific Beach (29), San Ysidro (25),
Test Zone (13), Spares (1). Unique areas: 42.

% ── 4. Methodology ────────────────────────────────────────────────
\section{Methodology}

\subsection{Data Preprocessing}

Several data quality issues were identified and addressed:

\begin{itemize}
    \item \textbf{Duplicates:} 12 duplicate transaction records were
    removed.

    \item \textbf{Invalid time ranges:} No rows had expiry before start
    time (0 invalid rows found).

    \item \textbf{Zero-duration transactions:} 44,806 records (2.29\%
    of total) had identical start and expiry timestamps. These were
    predominantly cash transactions and represent data recording errors.
    They were removed.

    \item \textbf{Test and dummy entries:} Locations labelled
    ``Test Zone'' (13 meters) and ``Spares'' (1 meter) were removed,
    including one entry with GPS coordinates of (0, 0) --- a clear
    placeholder entry.

    \item \textbf{Missing values:} Fields for \texttt{time\_limit},
    \texttt{days\_in\_operation}, and \texttt{price} had null values,
    filled with the string ``Unknown'' to preserve the rows.

    \item \textbf{Dataset joining:} The two datasets were joined using
    an inner join on the meter ID field (pole\_id in transactions,
    pole in locations). 3,334 meters were common to both datasets.
    The merged dataset contains 1,474,238 records.
\end{itemize}

\subsection{Feature Engineering}

Nineteen features were engineered from the raw data, grouped into
five categories as shown in Table~\ref{tab:features}.

\begin{table}[H]
\centering
\caption{Engineered Features by Category}
\label{tab:features}
\begin{tabular}{@{}ll@{}}
\toprule
\textbf{Category} & \textbf{Features} \\
\midrule
Time-based       & hour, day\_of\_week, month, week\_of\_year \\
Business flags   & is\_weekend, is\_morning\_peak (8--10am), \\
                 & is\_midday (11am--2pm), is\_afternoon (3--5pm) \\
Cyclical encoding & hour\_sin, hour\_cos, dow\_sin, dow\_cos \\
Transaction      & trans\_amt\_dollars \\
Location         & price\_per\_hr, time\_limit\_mins, operates\_weekend, \\
                 & zone\_enc, area\_enc, pay\_method\_enc \\
\bottomrule
\end{tabular}
\end{table}

Cyclical encoding was applied to hour and day-of-week features using
sine and cosine transformations, ensuring that adjacent time periods
are treated as similar (e.g. 11pm and midnight). Transaction duration
was deliberately excluded from the feature set to prevent data leakage,
as it directly determines the target variable.

\subsection{Target Variable}

The target variable \texttt{is\_long\_stay} is defined as:

\begin{equation}
\texttt{is\_long\_stay} =
\begin{cases}
1 & \text{if duration} > 60 \text{ minutes (Long Stay)} \\
0 & \text{if duration} \leq 60 \text{ minutes (Short Stay)}
\end{cases}
\end{equation}

Class distribution in the final dataset: Long Stay: 786,859 (53.4\%),
Short Stay: 687,379 (46.6\%). The dataset is slightly imbalanced,
which was addressed through class weighting in both models.

\subsection{Model Selection}

An 80/20 stratified train/test split was applied:
Training rows: 1,179,390. Test rows: 294,848.

\textbf{Random Forest:} An ensemble of 100 decision trees with
\texttt{max\_depth=12} and \texttt{class\_weight=`balanced'}.

\textbf{XGBoost:} A gradient boosting classifier with 100 estimators,
\texttt{max\_depth=6}, \texttt{learning\_rate=0.1}, and
\texttt{scale\_pos\_weight} set to the ratio of negative to positive
class samples.

% ── 5. Experimental Results ───────────────────────────────────────
\section{Experimental Results}

\subsection{Exploratory Data Analysis}

\begin{figure}[H]
    \centering
    \includegraphics[width=0.95\textwidth]{eda_charts.png}
    \caption{EDA charts --- Long stay rate by hour, day of week,
    zone, payment method, duration distribution, and correlation heatmap.}
    \label{fig:eda}
\end{figure}

Analysis of the cleaned dataset revealed the following patterns:

\begin{itemize}
    \item \textbf{Hour of day:} Long stay rate is highest at hour 6
    (100\%), meaning virtually every early-morning transaction is a
    long stay. The rate stabilizes around 50--55\% during core
    business hours (10am--4pm), dips sharply to 18\% at 5pm,
    then rises again to 60\% at 6pm.

    \item \textbf{Day of week:} All weekdays show long stay rates
    around 50--55\%. Saturday has the highest rate at approximately
    58\%. Sunday has almost no valid transactions since most San Diego
    meters do not operate on Sundays.

    \item \textbf{Zone:} City zone has the highest long stay rate
    ($\approx$70\%), followed by Downtown ($\approx$62\%), Uptown
    ($\approx$40\%), and Mid-City ($\approx$25\%).

    \item \textbf{Duration distribution:} The distribution is bimodal
    --- a large cluster of short transactions under 60 minutes, and a
    prominent spike at exactly 120 minutes (the most common time
    limit). This suggests many drivers pay for the maximum allowed
    time upfront.

    \item \textbf{Payment method:} Credit card payments have the
    highest long stay rate ($\approx$80\%), followed by Smart Card
    ($\approx$60\%) and Cash ($\approx$28\%).

    \item \textbf{Correlation heatmap:} \texttt{trans\_amt\_dollars}
    shows the strongest positive correlation with
    \texttt{is\_long\_stay} (0.63), confirming it as the dominant
    predictive signal. \texttt{hour} shows a negative correlation
    (-0.18), meaning later hours tend toward shorter stays.
\end{itemize}

\subsection{Model Performance}

Both models were evaluated on 294,848 test records. Results are
presented in Table~\ref{tab:results}.

\begin{table}[H]
\centering
\caption{Model Performance Comparison}
\label{tab:results}
\begin{tabular}{@{}lcccc@{}}
\toprule
\textbf{Model} & \textbf{Accuracy} & \textbf{Precision} &
\textbf{Recall} & \textbf{F1 Score} \\
\midrule
Random Forest & 0.8612 & 0.9387 & 0.7917 & 0.8590 \\
XGBoost       & 0.8640 & 0.9579 & 0.7796 & 0.8596 \\
\bottomrule
\end{tabular}
\end{table}

\subsection{Confusion Matrices}

\begin{figure}[H]
    \centering
    \includegraphics[width=0.95\textwidth]{confusion_matrices.png}
    \caption{Confusion matrices for Random Forest (left) and
    XGBoost (right).}
    \label{fig:confusion}
\end{figure}

\begin{table}[H]
\centering
\caption{Confusion Matrix Values (Test Set: 294,848 records)}
\label{tab:confusion}
\begin{tabular}{@{}lcccc@{}}
\toprule
\textbf{Model} & \textbf{True Short Stay} & \textbf{False Long Stay} &
\textbf{False Short Stay} & \textbf{True Long Stay} \\
\midrule
Random Forest & 129,342 & 8,134 & 32,781 & 124,591 \\
XGBoost       & 132,078 & 5,398 & 34,692 & 122,680 \\
\bottomrule
\end{tabular}
\end{table}

XGBoost makes significantly fewer false positives (5,398 vs 8,134),
explaining its higher precision (0.9579 vs 0.9387). Random Forest
catches more actual Long Stay cases (124,591 vs 122,680), giving it
slightly higher recall (0.7917 vs 0.7796).

\subsection{Feature Importance}

\begin{figure}[H]
    \centering
    \includegraphics[width=0.95\textwidth]{feature_importance.png}
    \caption{Feature importances for Random Forest (left) and
    XGBoost (right).}
    \label{fig:features}
\end{figure}

\textbf{Random Forest top features:}
\begin{enumerate}
    \item \texttt{trans\_amt\_dollars} (0.55) --- amount paid
    \item \texttt{pay\_method\_enc} (0.22) --- payment method
    \item \texttt{time\_limit\_mins} --- maximum allowed duration
    \item \texttt{hour} --- hour of day
    \item \texttt{zone\_enc} --- parking zone
\end{enumerate}

\textbf{XGBoost top features:}
\begin{enumerate}
    \item \texttt{trans\_amt\_dollars} (0.73) --- amount paid
    \item \texttt{zone\_enc} --- parking zone
    \item \texttt{time\_limit\_mins} --- maximum allowed duration
    \item \texttt{price\_per\_hr} --- hourly meter rate
    \item \texttt{hour\_sin} --- cyclical hour encoding
\end{enumerate}

\texttt{trans\_amt\_dollars} dominates both models, accounting for
55\% of Random Forest importance and 73\% of XGBoost importance.
This makes intuitive sense --- a driver paying for a larger block of
time upfront is a strong signal of intended long-stay parking.

\subsection{Prediction Function}

A prediction function was built to allow real-world inference:

\begin{verbatim}
predict_stay(hour=9, day_of_week=1, month=6, week_of_year=24,
             trans_amt_dollars=2.50, price_per_hr=2.50,
             time_limit_mins=120, operates_weekend=1,
             zone_name='Downtown', area_name='Little Italy',
             pay_method_name='CREDIT CARD')
\end{verbatim}

Output:
\begin{verbatim}
Random Forest -> Long Stay (probability: 0.95)
XGBoost       -> Long Stay (probability: 1.00)
\end{verbatim}

% ── 6. Discussion ─────────────────────────────────────────────────
\section{Discussion}

Both models perform well above baseline for this classification task.
The dominant influence of \texttt{trans\_amt\_dollars} is the most
notable finding --- a driver paying more upfront is a strong signal
of intended long-stay parking. This feature alone accounts for over
55\% of Random Forest importance and 73\% of XGBoost importance.

Payment method is the second most important feature in Random Forest,
suggesting that credit card users behave differently from cash payers
--- likely because card payments allow for larger, more convenient
transactions covering longer durations.

Time-of-day features (hour, hour\_sin) play a meaningful supporting
role. The sharp dip at 5pm (18\% long stay rate) is a notable pattern,
likely reflecting after-work visitors paying for just enough time to
complete a short errand before enforcement ends.

One limitation of the current approach is that the dataset only captures
paid transactions. Free parking events --- common in residential areas
and on Sundays --- are not represented. A more complete picture would
require sensor-based occupancy data alongside transaction records.

% ── 7. Conclusion ─────────────────────────────────────────────────
\section{Conclusion and Future Work}

This project demonstrated that real-world parking transaction data
can be used effectively to predict parking behavior with over 86\%
accuracy. XGBoost slightly outperforms Random Forest across most
metrics, and transaction amount is the dominant predictive feature
across both models.

Several directions for future work are worth exploring:

\begin{itemize}
    \item \textbf{Real-time integration:} Deploying the model as a
    REST API could allow integration with navigation applications,
    providing live parking guidance to drivers.

    \item \textbf{Weather and events data:} Adding external context
    (local events, weather) could improve predictions for atypical days.

    \item \textbf{Time series modeling:} Incorporating lag features
    (e.g. occupancy in the previous 30 minutes) could capture temporal
    dependencies not currently available in the feature set.

    \item \textbf{Multi-city generalization:} Testing the model on
    parking data from other cities would evaluate how well the learned
    patterns generalize across different urban contexts.
\end{itemize}

% ── References ────────────────────────────────────────────────────
\bibliographystyle{plain}

\begin{thebibliography}{9}

\bibitem{zheng2015}
Zheng, Y., Rajasegarar, S., \& Leckie, C. (2015).
\textit{Parking availability prediction for sensor-enabled car parks
in smart cities}.
In \textit{Proceedings of the IEEE 10th International Conference on
Intelligent Sensors, Sensor Networks and Information Processing}.
IEEE.

\bibitem{shoup2006}
Shoup, D. C. (2006).
\textit{Cruising for parking}.
\textit{Transport Policy}, 13(6), 479--486.

\bibitem{ji2014}
Ji, Y., Tang, D., Blythe, P., Guo, W., \& Wang, W. (2014).
\textit{Short-term forecasting of available parking space using
wavelet neural network model}.
\textit{IET Intelligent Transport Systems}, 9(2), 202--209.

\bibitem{chen2016}
Chen, T., \& Guestrin, C. (2016).
\textit{XGBoost: A scalable tree boosting system}.
In \textit{Proceedings of the 22nd ACM SIGKDD International
Conference on Knowledge Discovery and Data Mining} (pp. 785--794).

\bibitem{uci2016}
University of California, Irvine. (2016).
\textit{Parking Birmingham Data Set}.
UCI Machine Learning Repository.
\url{https://archive.ics.uci.edu/ml/datasets/Parking+Birmingham}

\bibitem{breiman2001}
Breiman, L. (2001).
\textit{Random forests}.
\textit{Machine Learning}, 45(1), 5--32.

\end{thebibliography}

\end{document}